\section{Conclusion}\label{sec:conclusion}

For linear elastostatics, the assembled discrete potential energy is an exact label-free
training objective: minimising it is supervised regression in the stiffness norm, which
measures the strain energy of the stress error, and it needs no solve. The same identity
makes the energy an evaluation metric that follows the stress rather than the displacement,
and a failure detector: the sign of the energy of a prediction tells whether it is worse
than the zero field, and a closed-form rescaling ensures that it is not. In
pre-registered comparisons on three-dimensional tetrahedral meshes, a transformer trained
on the energy alone was more accurate on every metric than the same transformer trained on
the solutions of the same instances, and a supervised graph network with lower
displacement errors had several times larger stress errors, so that displacement error
ranked the two in the opposite order to stress error. In two dimensions the label-free
network had a larger displacement error than the supervised one and smaller energy and stress
errors. Trained on the stiffness-norm error instead of the displacement error, with the same
labels, the supervised network halved its von Mises error, as pre-registered, and its errors
lay lower in the spectrum of the stiffness matrix, as measured afterwards: with the labels
fixed, the norm of the loss decided the accuracy of the stresses. In two dimensions, warm
starts from a label-free prediction did not shorten conjugate gradients to a tight
tolerance. In three dimensions no network transferred to a substantially finer mesh without
retraining, although a label-free rescaling removed most of the label-free network's error
there.
These limits, and the criteria that were not met, are reported with the results.
